% ID: SEL-D0043-OB-E03
% ESTADO: validada. S044, cotejada con examen y criterios.
Los 16 números tienen la misma probabilidad de ser elegidos.

\textbf{a)} Los impares son 225, 201, 193, 167, 215 y 171.
La probabilidad es $\boxed{6/16=3/8}$.

\textbf{b)} Los múltiplos de 5 son 225, 210, 180, 215, 120
y 190. Tres superan 200, así que la probabilidad es $\boxed{1/2}$.

\textbf{c)} Hay seis números mayores que 200 y diez pares.
Los que cumplen ambas condiciones son 210, 218 y 222.
Así, $P(S\cap T)=3/16$, distinto de
$P(S)P(T)=(6/16)(10/16)=15/64$. Son dependientes.

\textbf{d)} Por inclusión-exclusión, $P(S\cup T)=(6+10-3)/16$.
Resulta $\boxed{13/16}$.
